%==============================================================================
%  CAPÍTULO XVII
%==============================================================================
\chapter{Manual de tramitación: renegociación administrativa}
\label{cap:tram-renegociacion}

\fichaProcedimiento
  {Persona Deudora, incluidos quienes emiten boletas de honorarios.}
  {\Superir{}, por plataforma electrónica. No hay tribunal.}
  {De tres a seis meses según la agenda de audiencias.}
  {Gratuito para el deudor. El honorario del abogado, si lo hay, es el único costo.}
  {Acuerdo de renegociación o de ejecución; en su defecto, derivación a liquidación.}

El desarrollo dogmático se encuentra en el capítulo \ref{cap:renegociacion}.

%==============================================================================
\bloqueEncargo
%==============================================================================

Este procedimiento presenta una particularidad que condiciona el encargo: es gratuito y
está diseñado para ser tramitado por el propio deudor sin abogado. La pregunta legítima
es qué aporta el profesional.

La respuesta es concreta y debe explicitarse al cliente:

\begin{enumerate}
  \item \textbf{El control de admisibilidad.} La causa más frecuente de rechazo es
  documental. Una solicitud mal armada se declara inadmisible y el deudor pierde meses.

  \item \textbf{La determinación del pasivo.} La audiencia de determinación es
  contradictoria. El deudor sin asistencia acepta montos, intereses y recargos que
  habrían sido objetables.

  \item \textbf{El diseño de la propuesta.} Un plan de pagos que el deudor no puede
  cumplir produce un acuerdo incumplido, y con él el retorno al punto de partida con
  menos herramientas.

  \item \textbf{La detección de deudas de tratamiento especial.} El deudor no sabe que su
  CAE probablemente sobrevivirá, ni que la pensión de alimentos tiene régimen propio.
\end{enumerate}

\begin{alerta}[Sobre la admisibilidad y las ejecuciones notificadas]
La existencia de demandas ejecutivas notificadas incide en la admisibilidad. Este punto
debe verificarse en los registros del Poder Judicial \emph{antes} de ingresar la
solicitud, y no confiarse a la memoria del cliente, que frecuentemente no distingue entre
una carta de cobranza y una notificación judicial. La exclusión del \art{260} no alcanza,
en todo caso, a los juicios ejecutivos de origen laboral, que no obstan a la admisibilidad.
\end{alerta}

%==============================================================================
\bloqueEntrevista
%==============================================================================

\begin{cuestionario}[Admisibilidad]
  \pregunta{¿Cuántas deudas tiene vencidas, por qué monto, y hace cuánto tiempo?}
  {El \art{260} exige dos o más deudas vencidas por más de 90 días corridos, de
  obligaciones diversas, que sumen más de 80 UF.}

  \pregunta{¿Ha sido notificado personalmente por un receptor judicial de alguna
  demanda?}
  {Formulada así, la pregunta distingue la notificación real de la cobranza
  extrajudicial. Es el punto que decide la admisibilidad.}

  \pregunta{¿Ha iniciado antes un procedimiento concursal de cualquier tipo?}
  {Una resolución de término de liquidación firme impide una nueva liquidación voluntaria
  durante cinco años (\art{273 B}). Debe verificarse por RUT antes de definir la vía.}
\end{cuestionario}

\begin{cuestionario}[Capacidad de pago real]
  \pregunta{¿Cuál es su ingreso líquido mensual y cuáles sus gastos fijos
  indispensables?}
  {De la diferencia sale la única cifra que importa: cuánto puede pagar sin volver a
  incumplir.}

  \pregunta{¿Alguien depende económicamente de usted?}
  {Incide en el gasto indispensable y en la razonabilidad de la propuesta.}

  \pregunta{¿Su ingreso es estable o variable?}
  {En el deudor a honorarios, un plan de cuota fija sobre ingreso variable está destinado
  al incumplimiento.}
\end{cuestionario}

\begin{foco}[La cifra que define el procedimiento]
Ingreso líquido menos gasto indispensable, con margen de seguridad. Todo lo que se
proponga por sobre esa cifra es un acuerdo que se incumplirá. El límite legal del plan de
reembolso es un techo, no una meta.
\end{foco}

%==============================================================================
\bloqueDocumental
%==============================================================================

\begin{checklist}
  \item Cédula de identidad vigente.
  \item Informe consolidado de deudas y certificado de boletín comercial.
  \item Liquidaciones de remuneraciones o boletas de honorarios de los últimos meses.
  \item Certificado de cotizaciones previsionales.
  \item Comprobantes de gastos fijos: arriendo, servicios básicos, educación, salud.
  \item Certificado de anotaciones vigentes de vehículos y de dominio de inmuebles.
  \item Consulta de causas en el portal del Poder Judicial por RUT, para descartar
  ejecuciones notificadas.
  \item Documentación de codeudores solidarios y avales, conforme a las pautas de la
  \ncg{28} \citeyearpar{ncg-28}.
\end{checklist}

%==============================================================================
\bloqueFocos
%==============================================================================

\subsection{La audiencia de determinación del pasivo}

Es la audiencia decisiva y la que más se subestima. En ella se fija el monto sobre el que
se negociará. Debe llegarse con la nómina propia contrastada contra el informe
consolidado, y con objeciones preparadas respecto de intereses, comisiones y recargos
cuya procedencia no esté acreditada.

\subsection{La audiencia de renegociación}

Aquí se vota la propuesta. El resultado depende menos de la argumentación jurídica que de
la verosimilitud del plan. Un plan conservador y sostenible obtiene más apoyo que uno
ambicioso que los acreedores saben irrealizable.

\subsection{Los efectos de la resolución de admisibilidad}

La paralización de ejecuciones, la exclusión temporal de registros comerciales y la
suspensión de la prescripción son efectos inmediatos que el cliente debe conocer, porque
son la razón por la que vale la pena ingresar aun cuando la propuesta final sea modesta.

\begin{practica}[Los registros comerciales tras el acuerdo]
Existe jurisprudencia de protección que ha sancionado a instituciones financieras por
mantener registros negativos vigentes respecto de deudores con acuerdos firmes. Si el
cliente reporta que sigue apareciendo en registros después del acuerdo, esa vía está
disponible (capítulo \ref{cap:renegociacion}).
\end{practica}

%==============================================================================
\bloqueCronograma
%==============================================================================

\begin{checklist}[Secuencia de trabajo]
  \item \textbf{Antes de ingresar.} Verificación documental completa y consulta de causas
  judiciales por RUT.
  \item \textbf{Ingreso.} Solicitud por la plataforma de la \superir{}.
  \item \textbf{Examen de admisibilidad.} Respuesta rápida a los requerimientos de
  subsanación: el plazo es breve y su incumplimiento cierra el procedimiento.
  \item \textbf{Resolución de admisibilidad.} Se producen la suspensión de ejecuciones y
  de la prescripción.
  \item \textbf{Audiencia de determinación del pasivo.} Comparecencia con nómina propia y
  objeciones preparadas.
  \item \textbf{Audiencia de renegociación.} Votación de la propuesta.
  \item \textbf{Acuerdo o audiencia de ejecución.} Publicación del acuerdo aprobado, o
  paso a ejecución.
\end{checklist}

%==============================================================================
\bloqueErrores
%==============================================================================

\errorfrecuente{Ingresar sin verificar ejecuciones notificadas}
{La solicitud se declara inadmisible y el deudor pierde meses, además de la confianza en
el procedimiento.}
{Consultar las causas por RUT en el portal del Poder Judicial antes de ingresar.}

\errorfrecuente{Comparecer a la audiencia de determinación sin nómina propia}
{Se consolidan montos, intereses y comisiones que eran objetables, y la propuesta se
calcula sobre un pasivo inflado.}
{Llegar con la nómina contrastada contra el informe consolidado y las objeciones
redactadas.}

\errorfrecuente{Proponer el máximo que la ley permite}
{El acuerdo se incumple, y el deudor queda peor que antes.}
{Proponer sobre la capacidad real de pago, con margen de seguridad.}

\errorfrecuente{No advertir sobre el CAE ni sobre las pensiones de alimentos}
{El cliente cierra el procedimiento con las deudas que más le pesaban intactas.}
{Advertirlo por escrito en la primera reunión y consignarlo en el acta de diagnóstico.}

\errorfrecuente{Perder los plazos de subsanación de la plataforma}
{El procedimiento se cierra por inactividad.}
{Configurar alertas y revisar la plataforma con la misma disciplina con que se revisa el
\boletin{} en sede judicial.}
